\documentclass[11pt,a4paper]{article}

\usepackage[utf8]{inputenc}
\usepackage[T1]{fontenc}
\usepackage[polish]{babel}
\usepackage{geometry}
\usepackage{booktabs}
\usepackage{tabularray}
\usepackage{afterpage}
\usepackage{rotating}
\usepackage{longtable}
\usepackage{array}
\usepackage{xcolor}
\usepackage{hyperref}
\usepackage{titlesec}
\usepackage{parskip}
\usepackage{enumitem}
\usepackage{fancyhdr}
\usepackage{pdflscape}

\geometry{margin=2.5cm}

\definecolor{accent}{RGB}{30, 60, 100}

\titleformat{\section}{\large\bfseries\color{accent}}{}{0em}{}[\titlerule]
\titleformat{\subsection}{\normalsize\bfseries}{}{0em}{}
\titleformat{\subsubsection}{\normalsize\itshape}{}{0em}{}

\pagestyle{fancy}
\fancyhf{}
\rhead{\small Plan Zarządzania Komunikacją}
\lhead{\small Projekt: \textit{Feedback}}
\cfoot{\thepage}
\renewcommand{\headrulewidth}{0.4pt}

\hypersetup{colorlinks=true, linkcolor=accent, urlcolor=accent}

\begin{document}

% ── STRONA TYTUŁOWA ──────────────────────────────────────────────
\begin{titlepage}
\centering
\vspace*{2cm}
{\LARGE\bfseries\color{accent} Plan Zarządzania Komunikacją}\\[0.5cm]
{\large Projekt: \textit{Feedback}}\\[0.3cm]
{\normalsize Mechanika parowania jako klucz do sukcesu}\\[2cm]

\begin{tabular}{ll}
\textbf{Kierownik projektu:}   & Kacper Trzciński \\[0.3em]
\textbf{Przewodnicząca KS:}    & Aleksandra Sobótka \\[0.3em]
\textbf{Główny Dostawca:}      & Aleksander Płocharski \\[0.3em]
\textbf{Główny Użytkownik:}    & Tomasz Herman \\[0.3em]
\textbf{Zespół:}               & A. Bielicki, R. Głasek, J. Pączka, \\
                               & M. Popkowicz, F. Shramko, K. Trzciński \\[0.3em]
\textbf{Data opracowania:}     & \today \\[0.3em]
\end{tabular}

\vfill
{\small Zarządzanie Przedsięwzięciami Informatycznymi\\
Wydział Matematyki i Nauk Informacyjnych, Politechnika Warszawska}
\end{titlepage}

% ── SPIS TREŚCI ─────────────────────────────────────────────────
\tableofcontents
\newpage

% ════════════════════════════════════════════════════════════════
\section{Wstęp}
% ════════════════════════════════════════════════════════════════

\subsection{Przeznaczenie dokumentu}
Niniejszy dokument opisuje sposób organizacji komunikacji w projekcie \textit{Feedback} --- 
grze \textit{puzzle-action platformer} tworzonej przez 6-osobowy zespół studencki w ramach przedmiotu Zarządzanie Przedsięwzięciami Informatycznymi. 
Dokument określa, w jaki sposób, jak często i za pomocą jakich narzędzi będą wymieniać się informacjami wszyscy uczestnicy projektu.

\subsection{Po co zarządzamy komunikacją w projekcie?}
Projekt realizowany jest przez zespół pracujący równolegle nad kilkoma obszarami (backend, grafika, game design, testy). 
Bez ustrukturyzowanego przepływu informacji łatwo o rozbieżności między dokumentacją a stanem implementacji, 
pominięcie ważnych decyzji albo brak wiedzy o blokujących problemach. Zarządzanie komunikacją ma na celu:
\begin{itemize}[noitemsep]
    \item zapewnienie, że każdy członek zespołu zna swoje bieżące zadania i priorytety;
    \item umożliwienie szybkiego wykrywania i eskalowania ryzyk i zagadnień;
    \item utrzymanie regularnego kontaktu z Komitetem Sterującym i prowadzącymi przedmioty;
    \item stworzenie utrwalonych zapisów decyzji projektowych, które chronią projekt przy ewentualnej zmianie składu zespołu.
\end{itemize}

% ════════════════════════════════════════════════════════════════
\section{Procedury Komunikacji}
% ════════════════════════════════════════════════════════════════

\subsection{Formy komunikacji}

\begin{description}[style=nextline]
    \item[Spotkania tygodniowe (stand-upy)]
    Cotygodniowe, krótkie spotkania całego zespołu (do 30 minut). Każdy uczestnik odpowiada na trzy pytania: 
    co zrobił od ostatniego spotkania, co planuje zrobić do następnego, czy coś go blokuje.

    \item[Spotkania etapowe (przeglądy sprintu)]
    Na końcu każdego sprintu (ok. co 2--5 tygodni) odbywa się dłuższe spotkanie z udziałem Komitetu Sterującego. 
    Omawia się stan projektu, wyniki minionego sprintu, plan kolejnego oraz aktualne ryzyka.

    \item[Komunikacja asynchroniczna]
    Bieżąca wymiana informacji między członkami zespołu za pośrednictwem serwera Discord zespołu.
    Służy zadawaniu pytań, dzieleniu się postępem prac i szybkiemu rozwiązywaniu problemów niewymagających spotkania.

    \item[Komunikacja formalna]
    Komunikacja z prowadzącymi przedmioty i Komitetem Sterującym odbywa się za pośrednictwem serwera Discord specjalności --- przez dedykowany kanał lub wiadomości bezpośrednie. Wszelka korespondencja wychodząca poza zespół prowadzona jest przez kierownika projektu. Projekt realizowany jest w dwóch semestrach z różnymi prowadzącymi: Aleksandra Sobótka nadzoruje semestr dokumentacyjny (zimowy 2025/2026), natomiast Aleksander Płocharski --- semestr implementacyjny (letni 2026/2027). Komunikacja formalna kierowana jest każdorazowo do aktualnie prowadzącego.

    \item[Prezentacje]
    Prezentacja koncepcji projektu odbyła się 17.04.2026. Możliwa jest prezentacja
    podsumowująca semestr letni pod koniec maja lub na początku czerwca 2026. Prezentacja
    finalna z demo gry zaplanowana jest na początek semestru letniego 2026/2027 przed nowym
    rocznikiem.

    \item[Przeglądy kodu (\textit{code review})]
    Każda zmiana w repozytorium przechodzi przez mechanizm pull requestów --- co najmniej jedna osoba poza autorem przegląda i zatwierdza kod przed scaleniem.
\end{description}

\subsection{Zasady ogólne}

\begin{enumerate}[noitemsep]
    \item Decyzje podjęte ustnie są dokumentowane w formie notatki i udostępniane całemu zespołowi w ciągu 24 godzin.
    \item Problemy blokujące pracę są zgłaszane kierownikowi projektu niezwłocznie --- nie czekamy do najbliższego stand-upu.
    \item Wszelka komunikacja dotycząca zmian zakresu lub harmonogramu odbywa się przez kierownika projektu i jest dokumentowana pisemnie.
    \item Kierownik projektu jest jedynym przedstawicielem zespołu w komunikacji z prowadzącymi przedmioty i Komitetem Sterującym.
\end{enumerate}

% ════════════════════════════════════════════════════════════════
\section{Narzędzia i Techniki}
Narzędzia i techniki komunikacji przedstawiono w formie tabeli \ref{tab:narzedzia}. 
% ════════════════════════════════════════════════════════════════

% ════════════════════════════════════════════════════════════════
\section{Raportowanie}
% ════════════════════════════════════════════════════════════════

\subsection{Raport końca sprintu}

\begin{tabular}{ll}
\textbf{Autor:}         & Kierownik projektu (Kacper Trzciński) \\
\textbf{Odbiorcy:}      & Komitet Sterujący, prowadzący przedmiot \\
\textbf{Częstotliwość:} & Po każdym sprincie (ok. co 2--5 tygodni) \\
\textbf{Format:}        & Dokument PDF, dostarczany przez Discord \\
\end{tabular}

\vspace{0.5em}
Zawartość raportu:
\begin{itemize}[noitemsep]
    \item lista zadań zamkniętych w sprincie i porównanie z planem;
    \item aktualny stan harmonogramu i kamieni milowych;
    \item nowe ryzyka i zmiany w ocenie istniejących;
    \item zagadnienia wymagające decyzji Komitetu Sterującego;
    \item plan zadań na kolejny sprint.
\end{itemize}

\subsection{Notatka ze spotkania}

\begin{tabular}{ll}
\textbf{Autor:}         & Sekretarz spotkania (rotacyjnie, wyznaczany na jego początku) \\
\textbf{Odbiorcy:}      & Wszyscy uczestnicy spotkania \\
\textbf{Częstotliwość:} & Po każdym spotkaniu zespołu \\
\textbf{Format:}        & Dokument w Google Docs, udostępniany w ciągu 24h \\
\end{tabular}

\vspace{0.5em}
Zawartość notatki:
\begin{itemize}[noitemsep]
    \item data, lista uczestników, czas trwania;
    \item omówione tematy i podjęte decyzje;
    \item action items: zadanie --- osoba odpowiedzialna --- termin.
\end{itemize}

\subsection{Raport z testów}

\begin{tabular}{ll}
\textbf{Autor:}         & Wyznaczony tester z zespołu projektowego \\
\textbf{Odbiorcy:}      & Kierownik projektu, Komitet Sterujący, Główny Dostawca \\
\textbf{Częstotliwość:} & Jednorazowo, po zakończeniu fazy testów (marzec~2027) \\
\textbf{Format:}        & Dokument PDF \\
\end{tabular}

\vspace{0.5em}
Zawartość raportu:
\begin{itemize}[noitemsep]
    \item lista zgłoszonych błędów z podziałem na priorytety (krytyczne, wysokie, niskie);
    \item które błędy zostały naprawione, a które odłożone poza zakres;
    \item ocena stabilności wersji grywalnej;
    \item rekomendacja dotycząca gotowości do prezentacji finalnej.
\end{itemize}

% ════════════════════════════════════════════════════════════════
\section{Harmonogram Działań Komunikacyjnych}
Harmonogram działań komunikacyjnych przedstawiono w formie tabeli \ref{tab:harmonogram}.
% ════════════════════════════════════════════════════════════════



\vspace{1em}
\noindent\textbf{Przegląd dokumentu:} niniejszy plan podlega przeglądowi na końcu każdego etapu projektu. W razie istotnych zmian (zmiana składu zespołu, zmiana narzędzi) kierownik projektu aktualizuje dokument i informuje wszystkich interesariuszy.

% TABELE

\vspace*{\fill}
\clearpage

\begin{landscape}
    \begin{longtblr}[
        caption = {Harmonogram działań komunikacyjnych},
        label   = {tab:harmonogram},
    ]{
        colspec = {p{6cm} p{4cm} p{3.5cm} p{9cm}},
        rowhead = 1,
        hlines,
        row{1} = {font=\bfseries},
    }
    Zdarzenie & Kiedy & Uczestnicy & Forma / Kanał \\
    Stand-up tygodniowy
        & Co tydzień (ustalony dzień tygodnia)
        & Cały zespół
        & Discord (kanał głosowy) lub stacjonarnie, maks. 30 min \\
    Przegląd końca sprintu i planowanie kolejnego
        & Po każdym sprincie (co 4--6 tyg.)
        & Zespół + KS
        & Spotkanie głosowe lub stacjonarne, ok. 60 min \\
    Raport etapowy do KS i prowadzących
        & Po każdym sprincie
        & KP $\to$ KS, prowadzący
        & Discord (wiadomość bezpośrednia lub kanał) z załączonym PDF (raport + zaktualizowany harmonogram) \\
    Dostarczenie dokumentów do oceny
        & Wg kamieni milowych (patrz Karta projektu)
        & KP $\to$ prowadzący
        & Discord z załączonym PDF \\
    Code review (przegląd kodu)
        & Przy każdym pull requeście
        & Autor + co najmniej 1 recenzent
        & GitHub --- recenzja i zatwierdzenie przed scaleniem \\
    Zgłaszanie ryzyk i zagadnień
        & Na bieżąco (max. 48h od wykrycia)
        & Każdy członek $\to$ KP
        & Discord \texttt{\#ryzyka} lub bezpośrednia wiadomość do KP \\
    Przekazanie wersji testowej testerom
        & Po Feature Complete (23.02.2027)
        & Gł. Dostawca $\to$ testerzy
        & Link przez GitHub Releases + formularz zgłoszeń (Google Forms) \\
    Prezentacja koncepcji (pitch)
        & 17.04.2026
        & Zespół + prowadzący + KS
        & Prezentacja live (slajdy), 10 min \\
    Prezentacja finalna z demo gry
        & Pocz. semestru letniego 2026/2027
        & Zespół + nowy rocznik + prowadzący
        & Demo live + landing page projektu \\
    \end{longtblr}

    \clearpage
    \begin{longtblr}[
        caption = {Narzędzia i techniki komunikacji},
        label   = {tab:narzedzia},
    ]{
        width   = \linewidth,
        colspec = {>{\bfseries}X[1] X[1.5] X[4]},
        rowhead = 1,
        hlines,
        row{1} = {font=\bfseries},
    }
    Narzędzie & Zastosowanie & Szczegóły użycia \\
    Discord (serwer zespołu)
        & Codzienna komunikacja, stand-upy głosowe
        & Kanały tematyczne: \texttt{\#general}, \texttt{\#backend}, \texttt{\#grafika}, \texttt{\#dokumentacja}, \texttt{\#ryzyka}. Kanał \texttt{\#announce} zarezerwowany wyłącznie dla kierownika projektu --- tylko ogłoszenia i ważne decyzje. \\
    Discord (serwer specjalności)
        & Formalna komunikacja z prowadzącymi przedmiot i Komitetem Sterującym
        & Dedykowany kanał lub wiadomości bezpośrednie na serwerze specjalności. Wszelka korespondencja wychodząca poza zespół prowadzona przez kierownika projektu. Czas odpowiedzi: do 48h. \\
    GitHub (Issues + Milestones)
        & Zarządzanie zadaniami (Kanban), przeglądy kodu
        & Issues odzwierciedlają zadania z WBS i służą jako tablica Kanban (kolumny: To Do / In Progress / Done). Milestones odpowiadają sprintom. Pull requesty wymagają co najmniej jednej recenzji przed scaleniem. \\
    Google Drive + Google Docs
        & Dokumentacja projektowa, notatki ze spotkań, raporty
        & Wspólny folder projektu z podfolderami: \texttt{/Dokumenty}, \texttt{/Spotkania}, \texttt{/Raporty}. Prowadzący przedmioty mają dostęp tylko do odczytu. Preferowany format edycji dokumentów przed eksportem do PDF. \\
    ProjectLibre
        & Harmonogram projektu (diagram Gantta)
        & Plik \texttt{harmonogram.pod} przechowywany w Google Drive. Aktualizowany przez kierownika projektu po każdym sprincie; eksport do PDF dołączany do raportu etapowego. \\
    Google Forms + Sheets
        & Zbieranie i agregowanie feedbacku od testerów
        & Prosty formularz do zgłaszania błędów w fazie testów (luty--marzec 2027). Odpowiedzi automatycznie trafiają do arkusza Google Sheets, który służy jako rejestr błędów. \\
    \end{longtblr}
\end{landscape}

\end{document}